\documentclass[11pt,a4paper]{article}
\usepackage[T1]{fontenc}
\usepackage{lmodern}
\usepackage[margin=25mm]{geometry}
\usepackage{amsmath,amssymb,booktabs}
\usepackage{microtype}
\usepackage{listings}
\usepackage{xcolor}
\usepackage{xurl}
\usepackage[colorlinks=true,linkcolor=blue!50!black,urlcolor=blue!50!black]{hyperref}
\lstset{basicstyle=\small\ttfamily,columns=fullflexible,
  keepspaces=true,breaklines=true,showstringspaces=false,
  frame=single,rulecolor=\color{black!20},aboveskip=1em,belowskip=1em}
\newcommand{\code}[1]{\texttt{#1}}
\newcommand{\eq}{\mathrm{eq}}
\newcommand{\ordinary}{\mathrm{ordinary}}
\newcommand{\pair}{\mathrm{pair}}
\setlength{\emergencystretch}{2em}
\title{Grey heavy-flavour pair processes\\\large GRACE implementation report}
\author{}
\date{LaTeX edition: 30 September 2026}

\begin{document}
\maketitle
\noindent This document reproduces the explanation in
\path{doc/developer_guide/pair_processes.md}, with typeset equations and
a clarification of the effective-opacity convention. The build and test
record below is the existing record from 29 September 2026, not a new run.

\section{Scope and status}

This implementation adds selectable \code{equilibrium} and \code{evolved}
treatments; \code{legacy} remains the default. It treats \emph{neutrino pairs},
not creation of charged muons:
\begin{itemize}
\item electron--positron annihilation, $e^-+e^+\leftrightarrow\nu+\bar\nu$;
\item nucleon bremsstrahlung, $N+N\leftrightarrow N+N+\nu+\bar\nu$;
\item an explicitly optional transverse-plasmon model,
  $\gamma^*\leftrightarrow\nu+\bar\nu$.
\end{itemize}
These thermal channels do not require charged muons to be present. They must
not be multiplied by the charged-muon abundance/density/temperature gate.
Muonic nucleon charged-current reactions remain the responsibility of the
ordinary emission/absorption/scattering (EAS) provider. Muon decay/inverse
decay is a separate, optional extension documented in
\path{doc/developer_guide/muon_decay.rst}; it connects four species and changes
their source/backreaction workflow. The thermal-pair module described here
does not add muon--antimuon annihilation, flavour conversion, or missing
electron-flavour pair channels in the baseline.

This is an \emph{isotropic-kernel grey closure}, not an exact Boltzmann
solution. Conservation and detailed balance can be checked independently of
the accuracy of its spectral/angular assumptions. In particular, it is not a
validated polar-funnel annihilation/heating model and is not claimed to
outperform a spectral Monte Carlo calculation. Production use requires
quadrature, resolution, timestep and target-GPU validation.

\section{Configuration}

Build with 3 or 5 neutrino species and initialize the repository's existing
\path{extern/bns_nurates} submodule. Its pinned revision is
\path{efbabe0d4dfa6935c9ecbbd59a30c69b7cd1b91e}; no new revision is selected here.

\begin{lstlisting}
m1:
  eas:
    kind: neutrino_weakhub        # alternatively neutrino_analytic
    pair_treatment: evolved      # legacy | equilibrium | evolved
    pair_quadrature_order: 16     # 8..32; compare 16, 24, 32
    weakhub_pair_content: none    # ONLY after checking table provenance
    betaeq_policy: off
    pair_annihilation: true
    bremsstrahlung: true
    plasmon_decay: true
    pair_plasmon_model: transverse_mass
\end{lstlisting}

No spectral WeakHub tables are needed: the new pair kernels are evaluated
locally. The ordinary grey WeakHub table \emph{must exclude these
heavy-flavour thermal pair channels}. \path{weakhub_pair_content: none} is a
declaration, not a filter: the code cannot subtract unknown pair contributions
from an inclusive grey table. \code{unknown} and \code{included} are rejected
in the new modes. The old standalone \path{bns_nurates} EAS adapter is not
used by these modes.

The new modes reject algebraic \path{betaeq_policy} replacements of the
actual matter state. Ordinary temperature corrections remain confined to
ordinary rates. Pair rates receive neither an extra temperature-square
multiplier nor an optical-depth fugacity suppression. The historical muonic
$\eta=\pm5$ clamp is retained in legacy mode only: using it in ordinary
charged-current Kirchhoff factors but not in the new pair calculation would
introduce incompatible equilibrium targets. GRACE's existing equation-of-state
(EOS) chemical-potential convention is retained:
\[
 \mu_{\nu_\mu}=\mu_\mu+\mu_p-\mu_n-Q_{np},\qquad
 \mu_{\bar\nu_\mu}=-\mu_{\nu_\mu}.
\]

\section{Collision operator and grey equations}

In the matter frame, per physical helicity species, define
\[
 d\Pi=D(\epsilon)\,d\epsilon,\qquad
 D(\epsilon)=\frac{4\pi\epsilon^2}{(hc)^3},\qquad
 n=\int f\,d\Pi,\qquad J=\int\epsilon f\,d\Pi.
\]
Here $T$ is expressed in energy units. $R_a(\epsilon,\epsilon')$ is the
angle-independent annihilation kernel with the first argument belonging to
the neutrino, the second to the antineutrino. Kernel units are
$\mathrm{cm^3\,s^{-1}}$. For these thermal channels the matter bath has zero
net chemical potential for the event, so
\begin{equation}\label{eq:kernel-balance}
 R_p(\epsilon,\epsilon')
 =e^{-(\epsilon+\epsilon')/T}R_a(\epsilon,\epsilon').
\end{equation}
For a specified partner distribution,
\begin{equation}\label{eq:ja}
 j_\nu(\epsilon)=\int R_p(1-\bar f')\,d\Pi',\qquad
 a_\nu(\epsilon)=\int R_a\bar f'\,d\Pi',
\end{equation}
\begin{equation}\label{eq:collision}
 C_\nu=(1-f)j_\nu-fa_\nu
       =j_\nu-\lambda_\nu f,\qquad \lambda_\nu=j_\nu+a_\nu.
\end{equation}
The effective opacity in the rearranged source $j-\lambda f$ is
$\lambda/c$, not $a/c$. This does not introduce an additional physical
absorption process: $a$ is still the true inverse-reaction coefficient, and
$(1-f)j-af$ is an equally correct form. The extra $j$ in $\lambda$ accounts
for own-species Pauli blocking. Both $j$ and $a$ here have units
$\mathrm{s^{-1}}$, not the units of a grey energy emissivity.
The antineutrino equation uses the same event kernel with the roles of the
arguments transposed. Electron-pair kernels are not generally symmetric under
energy exchange.

For weights $w=1$ (number) or $w=\epsilon$ (energy), coefficients for a
specified spectral shape can be written
\begin{equation}\label{eq:grey}
 Q_w=\int wj\,d\Pi,\qquad
 c\kappa_w=\frac{\int w\lambda f\,d\Pi}{\int wf\,d\Pi}.
\end{equation}
Here $Q$ is spontaneous emission \emph{with partner blocking}, and
own-species blocking is contained in $\lambda$. Thus both emissivity and
opacity are computed from the kernel. Calling both of them independent
physical inputs would also be misleading: they obey microscopic detailed
balance, Eq.~\eqref{eq:kernel-balance}.

\subsection{Equilibrium partner}

Use $f^\eq=[\exp((\epsilon-\mu)/T)+1]^{-1}$ and the partner with $-\mu$.
Then
\begin{equation}\label{eq:kirchhoff}
 j(\epsilon)=\lambda(\epsilon)f^\eq(\epsilon),\qquad
 B_w=\int wf^\eq\,d\Pi,\qquad Q_w=c\kappa_w B_w,
\end{equation}
where the grey opacity in Eq.~\eqref{eq:grey} is weighted with $f^\eq$.
The code evaluates both sides from the same quadrature.
Equation~\eqref{eq:kirchhoff} is valid here because the blocking and averaging
conventions agree. It does not justify dividing an arbitrary unblocked
emissivity by a degenerate Fermi--Dirac (FD) blackbody.

The independent grey sources are $S_N=Q_N-c\kappa_N n$ and
$S_E=Q_E-c\kappa_E J$. They vanish at the prescribed local thermodynamic
equilibrium (LTE) target but \emph{need not give equal neutrino and
antineutrino number changes away from LTE}. This is an intentional limitation
of the equilibrium-partner comparison mode.

\subsection{Evolved partner}

At each trial of the implicit solve, GRACE's M1 closure supplies $J$,
$H^\alpha$ and $\Gamma$. The comoving number density is
$n=N_{\mathrm{lab}}/\Gamma$, with
$\Gamma=W(E-v^iF_i)/J$ in the existing grey number-current closure.
Reconstruct:
\begin{equation}\label{eq:reconstruction}
 \begin{gathered}
 f_i=[\exp(a+b\epsilon_i/\langle\epsilon\rangle)+1]^{-1},\qquad b>0,\\
 \sum_i W_i f_i=n,\qquad \sum_i W_i\epsilon_i f_i=J.
 \end{gathered}
\end{equation}
The fit parameters $a,b$ in Eq.~\eqref{eq:reconstruction} are unrelated to
the absorption coefficient $a_\nu$ in Eq.~\eqref{eq:ja}.
Both moments are matched on the \emph{same} positive quadrature that
evaluates the collisions, separately for each partner. $W_i$ includes
$D(\epsilon_i)$. This is a two-parameter FD/maximum-entropy spectral closure,
not a claim that the evolved spectrum is actually thermal. Invalid or
unrepresentable moments are rejected; the solver does not secretly replace
the partner by matter LTE.

For every ordered pair of energy nodes, accumulate one event rate:
\begin{equation}\label{eq:event}
 I_{ij}=W_iW_j\left\{R^p_{ij}(1-f_i)(1-\bar f_j)
                         -R^a_{ij}f_i\bar f_j\right\},
\end{equation}
\begin{equation}\label{eq:pair-moments}
 S_N=\bar S_N=\sum_{ij}I_{ij},\qquad
 S_E=\sum_{ij}\epsilon_iI_{ij},\qquad
 \bar S_E=\sum_{ij}\epsilon_jI_{ij}.
\end{equation}
Consequently the neutrino-number minus antineutrino-number source is zero
for pair-only interactions. The two energies need not be equal. The whole FD
family with common matter temperature and opposite chemical potentials is a
pair-equilibrium family; pair interactions alone do not select beta
equilibrium or erase a preexisting lepton-number difference. There is no
matter-Kirchhoff reset of Eqs.~\eqref{eq:event}--\eqref{eq:pair-moments}.

The grey flux spectrum is assumed proportional to the energy spectrum. With
$\lambda_E=\int\epsilon\lambda f\,d\Pi/J$, the covariant pair source is
\begin{equation}\label{eq:four-force}
 G^\alpha=S_Eu^\alpha-\lambda_EH^\alpha.
\end{equation}
In GRACE's undensitized Eulerian variables this contributes
\begin{equation}\label{eq:eulerian}
 \begin{aligned}
 \dot E&=\alpha W\left\{S_E-\lambda_E(E-v^iF_i-J)\right\},\\
 \dot F_i&=\alpha(Wv_iS_E-\lambda_EH_i),\qquad
 \dot N=\alpha S_N.
 \end{aligned}
\end{equation}
Ordinary charged-current/emission/scattering sources are added to
Eq.~\eqref{eq:eulerian} inside one coupled ten-variable diagonal
implicit--explicit (IMEX) solve:
\begin{equation}\label{eq:residual}
 \begin{gathered}
 U-U_*-h\,[S_\ordinary(U)+S_\pair(U)]=0,\\
 U=(E,F_x,F_y,F_z,N,\bar E,\bar F_x,\bar F_y,\bar F_z,\bar N),\qquad
 h=dt\,dtfact.
 \end{gathered}
\end{equation}
Matter and ordinary coefficients are held fixed over this local stage, as in
the existing GRACE source treatment. This is not a simultaneous EOS+matter
Newton solve. The existing conservative backreaction deposits the accepted
radiation energy/momentum and lepton-number changes into the fluid. Existing
backreaction limiters may throttle/reject those changes.

The solver uses scaled finite-difference Newton steps, positivity/cone checks
and line search. A continuation retry solves smaller coefficients of the same
original residual and ends at the full $h$; it is not a sequence of hidden
time substeps. Failure aborts explicitly, instead of accepting an unrelated
linear fallback. A final projection makes the two solved number equations
share a common pair increment; ordinary number absorption remains implicit.
The full nonlinear residual is checked again after this projection.
\path{pair_residual} records the normalized residual of the accepted state.

\subsection{Aggregate species and units}

Five species: indices 2/3 are separate $\nu_\mu/\bar\nu_\mu$; index 4
represents $\nu_\tau+\bar\nu_\tau$ with multiplicity 2. Three species: index 2
has multiplicity 4. Divide aggregate $E,F,N$ and ordinary emissivities by that
multiplicity before reconstructing a one-helicity spectrum. For a
charge-symmetric aggregate, use $[R(\epsilon,\epsilon')+
R(\epsilon',\epsilon)]/2$, evolve one representative spectrum and multiply
the output by the multiplicity. Do \emph{not} multiply its opacity by the
multiplicity.

Physical sources are converted using the same constants as existing EAS:
\[
 \begin{aligned}
 J_{\mathrm{code}}&=C_EJ,& n_{\mathrm{code}}&=C_Nn,\\
 Q_{E,\mathrm{code}}&=C_EQ_E/\code{TIMEGF},&
 Q_{N,\mathrm{code}}&=C_NQ_N/\code{TIMEGF},\\
 \lambda_{\mathrm{code}}&=\lambda/\code{TIMEGF}.&&
 \end{aligned}
\]
Here $C_E=\code{mev\_to\_erg}\,\code{RHOGF}\,\code{EPSGF}$ and
$C_N=\code{mnuc\_cgs}\,\code{RHOGF}$. Conserved variables are divided by
$\sqrt{\gamma}$ before solving, multiplied back afterwards.

\section{Kernel implementation and approximations}

Electron pairs use the pinned BNS\_NURATES Pons/Bruenn \code{PairPsi}
kernel, including its relativistic/massless-electron approximation (not an
exact finite-electron-mass kernel at arbitrarily low temperature),
algebraically evaluating its absorption coefficient without the unstable
product $\exp[(\epsilon+\epsilon')/T]R_p$. In the library's symbols the
implemented expression is
\[
 \begin{gathered}
 R_a^{ee}=-\frac{CT^2}{2[1-e^{-(y+z)}]}
 \left[\alpha_1^2\Psi_0(y,z,\eta_e)+
       \alpha_2^2\Psi_0(z,y,\eta_e)\right],\\
 y=\epsilon/T,\qquad z=\epsilon'/T,
 \end{gathered}
\]
with $C=\code{kBS\_Pair\_Phi}$ and the heavy-flavour weak couplings from
the pinned library; the result is converted from $\mathrm{nm^3/s}$ to
$\mathrm{cm^3/s}$. NN pairs use its unchanged Hannestad--Raffelt 1998 kernel,
without the optional Fischer medium correction. BNS\_NURATES uses nm units:
the numerical baryon density in $\mathrm{nm^{-3}}$ is its value in
$\mathrm{cm^{-3}}$ multiplied by $10^{-21}$, and the numerical kernel in
$\mathrm{cm^3/s}$ is its value in $\mathrm{nm^3/s}$ multiplied by $10^{-21}$.
The electron-pair expression retains the library's numerical
nonnegative-rate guard; this is not an independent revalidation of the
library's fit over every EOS state. Extremely degenerate states require
explicit convergence checks.

The optional plasma approximation uses two transverse polarisations, residue
$Z=1$, $\omega^2=k^2+m_\gamma^2$,
\begin{equation}\label{eq:plasmon-mass}
 \begin{gathered}
 m_\gamma^2=\frac{4\alpha_{\mathrm{EM}}}{3\pi}
 \left(\mu_e^2+\pi^2T^2/3\right),\\
 \Gamma_T(\omega)=\frac{G_F^2C_V^2m_\gamma^6}
 {48\pi^2\alpha_{\mathrm{EM}}\hbar\omega},\qquad
 C_V=-\tfrac12+2\sin^2\theta_W.
 \end{gathered}
\end{equation}
For $\omega=\epsilon+\epsilon'$ and $k^2=\omega^2-m_\gamma^2$, the allowed
region is $(\epsilon-\epsilon')^2\leq k^2$, with the transverse decay
distribution $dP/d\epsilon=3[1+(\epsilon-\epsilon')^2/k^2]/(4k)$.
Transforming the thermal photon phase space gives the joint event kernel
\begin{equation}\label{eq:plasmon-kernel}
 D(\epsilon)D(\epsilon')R_a=
 2\frac{4\pi}{(hc)^3}\,\omega\Gamma_T(\omega)
 [1+f_{\mathrm{BE}}(\omega)]\,\frac34
 \left[1+\frac{(\epsilon-\epsilon')^2}{k^2}\right].
\end{equation}
$R_p=\exp(-\omega/T)R_a$ includes $f_{\mathrm{BE}}$; the inverse rate
includes $1+f_{\mathrm{BE}}$, not $1-f_{\mathrm{BE}}$. This approximation
neglects longitudinal/axial/mixed modes, charged-muon polarisation, finite
electron mass corrections and the momentum-dependent transverse dispersion.
It is not quantitatively reliable for all cool/degenerate states; opt-in is
mandatory. Plasma thresholds can converge more slowly than smooth ee/NN
kernels on this tensor quadrature.

\section{EAS and evolution workflows}

These steps describe \code{muon\_decay: false}; see the separate decay guide
for the connected four-species extension.

\subsection{Equilibrium mode}
\begin{enumerate}
\item Read EOS and ordinary analytic/WeakHub rates; suppress old thermal extras.
\item Use the consistent muonic chemical potentials in ordinary Kirchhoff
  factors; apply the existing \emph{ordinary} gate and temperature corrections.
\item Construct thermal pair kernels and opposite-chemical-potential LTE spectra.
\item Compute both grey $Q$ and stimulated-absorption $\kappa$ from
  Eqs.~\eqref{eq:ja}--\eqref{eq:kirchhoff}.
\item Add these pair terms after ordinary corrections; preserve scattering.
\item Use the existing per-species $E,F,N$ implicit update and fluid backreaction.
\end{enumerate}

\subsection{Evolved mode}
\begin{enumerate}
\item Perform the same ordinary EAS calculation, excluding old thermal extras.
\item Store physical $T,\mu_e,n_b,X_n,X_p$ and enabled channels in pair
  auxiliaries. Ordinary EAS slots deliberately contain \emph{no pair terms}.
  Store current-state pair damping separately in \code{pair\_transport\_kappa*};
  flux corrections and optical-depth estimates use ordinary plus pair damping,
  while the collision solve reads only ordinary EAS coefficients.
\item In each implicit stage construct/cache one kernel matrix per physical pair.
\item Reconstruct both spectra from each trial $J,N/\Gamma$; compute
  Eqs.~\eqref{eq:event}--\eqref{eq:eulerian}.
\item Solve Eq.~\eqref{eq:residual}, including ordinary EAS, simultaneously
  for the two partners. Treat the charge-symmetric residual heavy aggregate
  with its multiplicity.
\item Run unchanged electron/photon source updates, then existing backreaction.
\end{enumerate}

Equilibrium-blended initial data temporarily use equilibrium pair EAS even
when the simulation selects evolved partners. The usual final EAS pass then
restores ordinary-only slots and material snapshots for evolution. Atmosphere
cells disable pair channels and clear their material/residual auxiliaries.

\section{Code map, review and validation}

\begin{itemize}
\item \path{neutrino_pair_collision.hh}: quadrature, kernels, FD
  reconstruction, common blocked collision integrals and units.
\item \path{neutrino_pair_update.hh}: coupled M1 residual and local nonlinear
  solver.
\item \path{eas_kinds.hh}, \path{parameters/m1.yaml}: setup selectors and
  incompatible-mode checks.
\item \path{eas_policies.hh}: ordinary/pair separation, matching muonic LTE
  targets, equilibrium coefficients and evolved material snapshots.
\item \path{variable_indices.hh/.cpp}: material/channel/residual auxiliary
  registration.
\item \path{m1.hh}: one-cell species mapping, densitization and aggregate factors.
\item \path{evolve.cpp}: pair-stage launch before ordinary sources/backreaction.
\item \path{m1.cpp}: setup validation, small-tile EAS launch and initialization
  support.
\item \path{CMakeLists.txt}, \path{test/CMakeLists.txt},
  \path{test_m1_pairs.cpp}: headers and tests. The Kokkos test helper now
  propagates the bundled YAML/p4est/HDF5 dependencies; existing backreaction
  tests support both Cowling and Z4 metric layouts. Two legacy-rate assertions
  were corrected to test the already-existing smooth muon suppression factor,
  rather than assuming a hard cutoff. The legacy rate implementation itself
  was not changed for those assertions.
\end{itemize}

The hot path has fixed-size automatic storage and Kokkos device functions, no
host-only solver, dynamic allocation or CUDA/HIP-specific API. Kernel matrices
are computed once per stage, not once per Newton trial. Separate small tiles
avoid the ordinary kernel's GPU launch-bound assumptions. This is portable
source design, not a GPU performance measurement; register pressure, spills
and throughput still need measurement on each target backend.

Tests cover unit normalization, FD moments, kernel orientation, LTE balance,
blocking/vacuum emission, unequal populations, implicit stiffness,
simultaneous ordinary EAS, moving-fluid LTE and execution on the configured
Kokkos backend. Run
\begin{lstlisting}
ctest --test-dir build-pairs-cpu -R m1_pair_test --output-on-failure
\end{lstlisting}
Do not equate unit-test success with a validated neutron-star evolution.

\subsection{Verified build (2026-09-29)}

The complete \code{grace} executable built with GCC/WSL, Kokkos Serial,
Debug, 3 spatial dimensions, 5 neutrino species, muons enabled and Cowling
metric. The following suites passed after the review/fix cycles:
\begin{center}
\begin{tabular}{lrr}
\toprule
Suite & Cases & Assertions \\
\midrule
\path{m1_pair_test} & 9 & 119 \\
\path{m1_analytic_rates_test} & 33 & 1655 \\
\path{m1_backreaction_test} & 44 & 643 \\
\midrule
Total & 86 & 2417 \\
\bottomrule
\end{tabular}
\end{center}
\begin{lstlisting}
cmake -S . -B build-pairs-cpu -DGRACE_USE_BUNDLED_DEPS=ON \
  -DGRACE_M1_NU_SPECIES=5 -DGRACE_3D=ON -DGRACE_METRIC_EVOL=COWLING \
  -DGRACE_PERF_TUNING=generic -DCMAKE_BUILD_TYPE=Debug
cmake --build build-pairs-cpu --target grace m1_pair_test \
  m1_analytic_rates_test m1_backreaction_test -j6
ctest --test-dir build-pairs-cpu \
  -R '^m1_(pair|analytic_rates|backreaction)_test$' --output-on-failure
\end{lstlisting}

Not verified here: GPU compilation/performance, a three-species executable,
full coupled hydrodynamic/merger evolutions, all extreme EOS states, or
quantitative accuracy against full angular/spectral transport. The standard
M1 flux cone is checked; this is not a full fermionic angular-realizability
closure. Even with matched chemical potentials, differences between finite
pair quadrature and the existing ordinary-EAS Fermi-integral fits leave a
quadrature-level LTE mismatch: check order convergence before production.
The existing generic implicit-solver counters describe the ordinary solver;
\path{pair_residual} is the new coupled-pair solver's diagnostic.

\section{Literature and what can actually be criticized}

\href{https://arxiv.org/html/2309.03526v2}{Ng et al., WeakHub}, give the pair
collision operator and reaction kernels (Sec.~3.4).
\href{https://arxiv.org/html/2412.04570v2}{Chiesa et al., BNS\_NURATES},
describe the portable kernel library and reconstruction from grey moments.
The new code reuses the kernels, not the library's complete grey closure.
\href{https://arxiv.org/abs/hep-ph/9302213}{Braaten and Segel} give the more
complete plasma treatment; the simplified model
\eqref{eq:plasmon-mass}--\eqref{eq:plasmon-kernel} must not be confused with it.

\href{https://arxiv.org/html/2604.14225v2}{Gieg et al.}, Appendix~A, drop
neutrino blocking in A4 before writing FD detailed balance in A5. That cannot
be the exact identity of the fully blocked collision operator; an additional
closure choice is involved. This implementation keeps both blocking factors
instead. \href{https://arxiv.org/html/2606.27425v1}{Foucart}, Sec.~II.4, uses
evolved spectral estimates with an opaque-region equilibrium fallback and an
angular correction; it is not simply an equilibrium-partner treatment. A
blanket claim that this new grey $R_0$ implementation is more accurate is
unwarranted. No unpublished criticism attributed to Harry is treated as
evidence of a specific error.

\end{document}
